\documentclass[12pt, a4paper]{book}

% ─── Core packages ────────────────────────────────────────────────────────────
\usepackage[T1]{fontenc}
\usepackage[utf8]{inputenc}
\usepackage[charter, cal=cmcal]{mathdesign}
\usepackage[scaled=0.92]{helvet}
\usepackage{geometry}
\geometry{
  a4paper,
  top=30mm, bottom=30mm,
  inner=32mm, outer=25mm,
  marginparwidth=18mm
}

% ─── Mathematics ──────────────────────────────────────────────────────────────
\usepackage{amsmath, amsthm}

% ─── Chemistry ────────────────────────────────────────────────────────────────
\usepackage[version=4]{mhchem}

% ─── Graphics & colour ────────────────────────────────────────────────────────
\usepackage{graphicx}
\usepackage[dvipsnames,svgnames,table]{xcolor}
\usepackage{tikz}
\usetikzlibrary{arrows.meta, shapes.geometric, positioning,
                calc, backgrounds, decorations.pathmorphing,
                decorations.text, patterns}
\usepackage{pgfplots}
\pgfplotsset{compat=1.18}
\usepackage{tcolorbox}
\tcbuselibrary{skins, breakable, theorems}

% ─── Layout ───────────────────────────────────────────────────────────────────
\usepackage{titlesec}
\usepackage{fancyhdr}
\usepackage{enumitem}
\usepackage{booktabs}
\usepackage{array}
\usepackage{multicol}
\usepackage{marginnote}
\usepackage{caption}
\usepackage{parskip}
\usepackage{pifont}
\usepackage{siunitx}
\usepackage[colorlinks=true, linkcolor=BrickRed,
            urlcolor=NavyBlue, citecolor=ForestGreen]{hyperref}

% ══════════════════════════════════════════════════════════════════════════════
% COLOUR PALETTE
% ══════════════════════════════════════════════════════════════════════════════
\definecolor{JourneyColor}{HTML}{065F46}   % forest green — journey headings
\definecolor{ModuleColor}{HTML}{0F766E}    % deep teal — module headings
\definecolor{PuzzleColor}{HTML}{B45309}    % deep amber — puzzle boxes
\definecolor{DiscoveryColor}{HTML}{D97706} % golden amber — discovery boxes
\definecolor{WarningColor}{HTML}{E11D48}   % coral red — "model breaks" boxes
\definecolor{ProblemColor}{HTML}{0F766E}   % dark teal — problem sets
\definecolor{ReflectColor}{HTML}{047857}   % green — reflection boxes
\definecolor{HistoryColor}{HTML}{9A3412}   % brown — history notes
\definecolor{LightBg}{HTML}{F8FAFC}        % near-white background
\definecolor{SynthColor}{HTML}{4A235A}

% ══════════════════════════════════════════════════════════════════════════════
% TCOLORBOX ENVIRONMENTS
% ══════════════════════════════════════════════════════════════════════════════
\newtcolorbox{puzzlebox}[1][]{
  enhanced, breakable,
  colback=PuzzleColor!6, colframe=PuzzleColor,
  fonttitle=\bfseries\sffamily, title={The Puzzle},
  attach boxed title to top left={yshift=-2mm, xshift=6mm},
  boxed title style={colback=PuzzleColor, colframe=PuzzleColor},
  borderline west={3pt}{0pt}{PuzzleColor}, #1}

\newtcolorbox{discoverybox}[1][]{
  enhanced, breakable,
  colback=DiscoveryColor!8, colframe=DiscoveryColor,
  fonttitle=\bfseries\sffamily, title={Discovery},
  attach boxed title to top left={yshift=-2mm, xshift=6mm},
  boxed title style={colback=DiscoveryColor, colframe=DiscoveryColor},
  borderline west={3pt}{0pt}{DiscoveryColor}, #1}

\newtcolorbox{problembox}[1][]{
  enhanced, breakable,
  colback=ProblemColor!5, colframe=ProblemColor,
  fonttitle=\bfseries\sffamily, title={Problems},
  attach boxed title to top left={yshift=-2mm, xshift=6mm},
  boxed title style={colback=ProblemColor, colframe=ProblemColor},
  borderline west={3pt}{0pt}{ProblemColor}, #1}

\newtcolorbox{deepbox}[1][]{
  enhanced, breakable,
  colback=JourneyColor!5, colframe=JourneyColor,
  fonttitle=\bfseries\sffamily, title={Deep Problems},
  attach boxed title to top left={yshift=-2mm, xshift=6mm},
  boxed title style={colback=JourneyColor, colframe=JourneyColor},
  borderline west={4pt}{0pt}{JourneyColor}, #1}

\newtcolorbox{historybox}[1][]{
  enhanced, breakable,
  colback=HistoryColor!7, colframe=HistoryColor!60,
  fonttitle=\bfseries\sffamily\small, title={Historical Note},
  attach boxed title to top left={yshift=-2mm, xshift=6mm},
  boxed title style={colback=HistoryColor!80, colframe=HistoryColor}, #1}

\newtcolorbox{reflectionbox}[1][]{
  enhanced, breakable,
  colback=ReflectColor!6, colframe=ReflectColor,
  fonttitle=\bfseries\sffamily, title={Reflection},
  attach boxed title to top left={yshift=-2mm, xshift=6mm},
  boxed title style={colback=ReflectColor, colframe=ReflectColor},
  borderline west={3pt}{0pt}{ReflectColor}, #1}

\newtcolorbox{synthbox}[1][]{
  enhanced, breakable,
  colback=SynthColor!5, colframe=SynthColor,
  fonttitle=\bfseries\sffamily, title={Journey Synthesis},
  attach boxed title to top left={yshift=-2mm, xshift=6mm},
  boxed title style={colback=SynthColor, colframe=SynthColor},
  borderline west={4pt}{0pt}{SynthColor}, #1}

\newtcolorbox{namedconcept}[2][]{
  enhanced,
  colback=ModuleColor!6, colframe=ModuleColor,
  fonttitle=\bfseries\sffamily\large, title={#2},
  attach boxed title to top left={yshift=-2mm, xshift=6mm},
  boxed title style={colback=ModuleColor, colframe=ModuleColor}, #1}

\newtcolorbox{openquestion}[1][]{
  enhanced, breakable,
  colback=DiscoveryColor!5, colframe=DiscoveryColor!70,
  fonttitle=\bfseries\sffamily, title={Open Question},
  attach boxed title to top left={yshift=-2mm, xshift=6mm},
  boxed title style={colback=DiscoveryColor!80, colframe=DiscoveryColor},
  borderline west={3pt}{0pt}{DiscoveryColor!70}, #1}

% ══════════════════════════════════════════════════════════════════════════════
% THEOREM ENVIRONMENTS
% ══════════════════════════════════════════════════════════════════════════════
\theoremstyle{definition}
\newtheorem{observation}{Observation}[chapter]
\newtheorem{principle}{Principle}[chapter]

% ══════════════════════════════════════════════════════════════════════════════
% SECTION FORMATTING
% ══════════════════════════════════════════════════════════════════════════════
\titleformat{\chapter}[display]
  {\normalfont\sffamily\huge\bfseries\color{JourneyColor}}
  {\textcolor{JourneyColor!60}{\chaptertitlename\ \thechapter}}
  {10pt}{\Huge}[\vspace{2pt}\textcolor{JourneyColor}{\titlerule[1.5pt]}]

\titleformat{\section}
  {\normalfont\sffamily\Large\bfseries\color{ModuleColor}}
  {\thesection}{1em}{}[\textcolor{ModuleColor!40}{\titlerule[0.6pt]}]

\titleformat{\subsection}
  {\normalfont\sffamily\large\bfseries\color{JourneyColor!80}}
  {\thesubsection}{1em}{}

\titleformat{\subsubsection}
  {\normalfont\sffamily\normalsize\bfseries}
  {\thesubsubsection}{1em}{}

% ══════════════════════════════════════════════════════════════════════════════
% HEADER / FOOTER
% ══════════════════════════════════════════════════════════════════════════════
\pagestyle{fancy}
\fancyhf{}
\fancyhead[LE]{\small\sffamily\textcolor{JourneyColor}{\leftmark}}
\fancyhead[RO]{\small\sffamily\textcolor{ModuleColor}{\rightmark}}
\fancyfoot[C]{\small\thepage}
\renewcommand{\headrulewidth}{0.4pt}

% ══════════════════════════════════════════════════════════════════════════════
% CUSTOM COMMANDS
% ══════════════════════════════════════════════════════════════════════════════
\newcommand{\chem}[1]{\ce{#1}}
\newcommand{\important}[1]{\textcolor{ModuleColor}{\textbf{#1}}}
\newcommand{\newterm}[1]{\textit{\textcolor{DiscoveryColor}{#1}}}
\newcommand{\sidenote}[1]{%
  \marginnote{\footnotesize\sffamily\textcolor{JourneyColor!70}{#1}}}

% ══════════════════════════════════════════════════════════════════════════════
% DOCUMENT
% ══════════════════════════════════════════════════════════════════════════════
\begin{document}
\setcounter{chapter}{0}
\chapter[Journey 1 --- What Is Stuff?]
        {Journey One\\[8pt]
         \textcolor{ModuleColor}{\large What Is Stuff?\ \ Matter and Its Patterns}}
\setcounter{section}{4}

% ══════════════════════════════════════════════════════════════════════════════
%  MODULE 1.5 — THE COMPLETE PICTURE: JOURNEY ONE SYNTHESIS
% ══════════════════════════════════════════════════════════════════════════════
\section{Module 1.5 --- The Complete Picture: Putting Journey One Together}
\label{sec:mod1.5}

\subsection*{What This Module Is}

This module is different from the previous four.
There are no new facts here.
Instead, we step back and look at the whole landscape
we have crossed, ask what the discoveries of Journey~1
actually mean as a \emph{way of doing science},
and identify the questions that Journey~1 answers ---
and, crucially, the deeper questions it opens.

The goal is synthesis: not just knowing the facts,
but understanding the structure of the argument.

% ─── Section: The Narrative Arc ───────────────────────────────────────────────
\subsection{The Narrative Arc of Journey One}

Journey~1 began with a question so old it seems trivial:
\emph{why do different substances behave differently?}
It ended with atoms.
The path between them was not a straight line.
It was a zigzag of observations, failed theories,
unexpected data, and forced conclusions.

Here is the complete arc:

\begin{center}
\begin{tikzpicture}[
  every node/.style={font=\sffamily\footnotesize, align=center},
  stepbox/.style={draw=JourneyColor!50, fill=JourneyColor!7,
                  rounded corners=4pt, minimum width=3.6cm,
                  minimum height=1.0cm, text width=3.3cm},
  failbox/.style={draw=WarningColor!60, fill=WarningColor!6,
                  rounded corners=4pt, minimum width=3.0cm,
                  minimum height=0.8cm, text width=2.8cm,
                  font=\sffamily\tiny, align=center},
  arr/.style={-{Stealth[length=5pt]}, thick, color=ModuleColor!70},
  failarr/.style={-{Stealth[length=4pt]}, thick, color=WarningColor!60,
                  dashed}
]

% Main chain - vertical spine
\node[stepbox] (obs)   at (0,  0)   {Observation:\\substances differ\\in properties};
\node[stepbox] (class) at (0, -2.2) {Classification:\\elements, compounds,\\mixtures};
\node[stepbox] (phlog) at (0, -4.4) {Burning and\\reactions observed:\\phlogiston theory};
\node[stepbox] (lav)   at (0, -6.6) {Lavoisier weighs\\everything:\\mass is conserved};
\node[stepbox] (proust)at (0, -8.8) {Fixed mass ratios\\in compounds:\\definite proportions};
\node[stepbox] (dalton)at (0,-11.0) {Multiple proportions:\\integer ratios demand\\discrete particles};
\node[stepbox] (avog)  at (0,-13.2) {Gas volumes confirm\\molecule formulas:\\Avogadro's hypothesis};
\node[stepbox, draw=DiscoveryColor, fill=DiscoveryColor!10,
               minimum height=1.3cm]
               (atom)  at (0,-15.6) {\textbf{Atomic-Molecular Theory}\\Matter = discrete atoms\\Each element = one atom type\\Compounds = fixed atom ratios};

% Arrows down the spine
\foreach \a/\b in {obs/class, class/phlog, phlog/lav, lav/proust,
                   proust/dalton, dalton/avog, avog/atom}
  \draw[arr] (\a) -- (\b);

% Failed theories and corrections on the right
\node[failbox] (phlogfail) at (5.5, -4.4)
  {Phlogiston theory:\\explains much,\\but mass increase\\breaks it};
\draw[failarr] (phlog.east) -- (phlogfail.west);
\draw[arr, color=DiscoveryColor!70] (phlogfail.south)
  -- ++(0,-0.6) -- ++(-2.6,-0.6) -- (lav.east);

\node[failbox] (daltonwrong) at (5.5, -11.0)
  {\ce{H2O} or \ce{HO}?\\Atomic masses\\ambiguous from\\mass data alone};
\draw[failarr] (dalton.east) -- (daltonwrong.west);
\draw[arr, color=DiscoveryColor!70] (daltonwrong.south)
  -- ++(0,-0.5) -- ++(-2.6,-0.9) -- (avog.east);

% Annotations on left
\node[font=\sffamily\tiny\itshape, text=HistoryColor,
      text width=2.5cm, align=right]
  at (-3.2, -1.1) {Module 1.1};
\node[font=\sffamily\tiny\itshape, text=HistoryColor,
      text width=2.5cm, align=right]
  at (-3.2, -3.3) {Module 1.1};
\node[font=\sffamily\tiny\itshape, text=HistoryColor,
      text width=2.5cm, align=right]
  at (-3.2, -5.5) {Module 1.2};
\node[font=\sffamily\tiny\itshape, text=HistoryColor,
      text width=2.5cm, align=right]
  at (-3.2, -7.7) {Module 1.2--1.3};
\node[font=\sffamily\tiny\itshape, text=HistoryColor,
      text width=2.5cm, align=right]
  at (-3.2, -9.9) {Module 1.3};
\node[font=\sffamily\tiny\itshape, text=HistoryColor,
      text width=2.5cm, align=right]
  at (-3.2, -12.1) {Module 1.3};
\node[font=\sffamily\tiny\itshape, text=HistoryColor,
      text width=2.5cm, align=right]
  at (-3.2, -14.3) {Module 1.4};

\end{tikzpicture}
\end{center}

\vspace{0.5em}

Notice the structure.
Twice in this journey, a theory that explained most of the data
was overturned by a \emph{quantitative anomaly} that the theory
could not absorb:
phlogiston failed because metals gained mass on burning,
and Dalton's simple-atom picture failed because gas volumes
demanded diatomic molecules.

In both cases, the anomaly was not immediately obvious.
It required careful measurement.
It required someone to insist on accounting for everything.

\emph{This is how science actually works.}
Not by grand theoretical leaps, but by small quantitative
inconsistencies that accumulate until a better framework
becomes unavoidable.

% ─── Section: Five Laws, One Picture ──────────────────────────────────────────
\subsection{Five Laws, One Picture}

Journey~1 produced five empirical laws.
Each one was discovered by different people,
in different countries, using different apparatus.
Yet they all say the same thing from different angles.

\begin{center}
\begin{tikzpicture}[
  every node/.style={font=\sffamily\small, align=center},
  lawbox/.style={draw=JourneyColor!60, fill=JourneyColor!8,
                 rounded corners=5pt, minimum width=5.5cm,
                 minimum height=1.1cm, text width=5.2cm},
  corebox/.style={draw=DiscoveryColor, fill=DiscoveryColor!12,
                  rounded corners=6pt, minimum width=4.8cm,
                  minimum height=2.8cm, text width=4.5cm},
  arr/.style={-{Stealth[length=5pt]}, thick, color=JourneyColor!60}
]
  % Five laws around the outside
  \node[lawbox] (cm) at (0,    5.0)
    {\textbf{Conservation of Mass}\\(Lavoisier, 1785)\\
     Mass in = Mass out};
  \node[lawbox] (dp) at (5.5,  2.2)
    {\textbf{Definite Proportions}\\(Proust, 1799)\\
     Fixed mass ratios in compounds};
  \node[lawbox] (mp) at (4.2, -2.5)
    {\textbf{Multiple Proportions}\\(Dalton, 1803)\\
     Ratios form small integers};
  \node[lawbox] (cv) at (-4.2,-2.5)
    {\textbf{Combining Volumes}\\(Gay-Lussac, 1809)\\
     Gas volumes in integer ratios};
  \node[lawbox] (av) at (-5.5, 2.2)
    {\textbf{Avogadro's Hypothesis}\\(1811)\\
     Equal volumes = equal molecules};

  % Central core
  \node[corebox] (core) at (0, 0.8)
    {\textbf{Atomic-Molecular\\Theory}\\[5pt]
     \small Matter is made of\\
     discrete, indivisible atoms.\\
     Compounds have fixed\\
     atom ratios.\\
     Equal gas volumes $\Leftrightarrow$\\
     equal molecule counts.};

  % Arrows from laws to core
  \foreach \law in {cm, dp, mp, cv, av}
    \draw[arr] (\law) -- (core);

\end{tikzpicture}
\end{center}

\vspace{0.5em}

The atomic-molecular theory is not one of these laws.
It is the \emph{explanation} of all five simultaneously.
When a single picture explains five independent sets
of measurements, that is strong evidence the picture
is pointing at something real.

\begin{principle}
\textbf{Convergence as evidence.}
When multiple independent lines of evidence all point
to the same underlying structure, the probability that
the structure is real increases dramatically.
No single law proves atoms exist.
But five laws that \emph{all follow from} the atomic picture ---
that is a compelling argument.
\end{principle}

% ─── Section: What the Theory Explains ────────────────────────────────────────
\subsection{Accounting for Everything: A Checklist}

Let us be precise about what the atomic-molecular theory
explains and what it does not.

\begin{center}
\begin{tabular}{p{6cm} p{2.2cm} p{5.5cm}}
\toprule
\textbf{Phenomenon} & \textbf{Explained?} & \textbf{How} \\
\midrule
Mass conserved in reactions
  & \checkmark
  & Atoms rearrange; none created/destroyed \\[4pt]
Fixed mass ratios in compounds
  & \checkmark
  & Each compound has a fixed atom ratio \\[4pt]
Integer ratios between related compounds
  & \checkmark
  & Atoms are discrete; you cannot have half an atom \\[4pt]
Gas volume ratios are integers
  & \checkmark
  & Equal volumes contain equal numbers of molecules \\[4pt]
\emph{Why} elements have characteristic masses
  & \ding{55}
  & Theory assumes it; no explanation \\[4pt]
\emph{Why} atoms of different elements differ
  & \ding{55}
  & Internal structure of atoms unknown \\[4pt]
\emph{Why} certain atoms bond together
  & \ding{55}
  & Bonding mechanism not yet known \\[4pt]
Periodic patterns in element properties
  & \ding{55}
  & Requires knowledge of atomic structure \\[4pt]
Why matter has the properties it has
  & \ding{55}
  & Deeper structure still hidden \\
\bottomrule
\end{tabular}
\end{center}

\vspace{0.5em}

The right-hand column of failures is not a weakness.
It is a map of where we are going next.

Every \ding{55} in the table is a door.
Journey~2 (architecture of the atom) will open several of them.
Journey~3 (the periodic table) will open more.
By Journey~4 (the chemical bond), almost all of them
will have answers.

\begin{puzzlebox}
Each unanswered question in the table above
led to a major scientific revolution.

\begin{enumerate}[label=(\alph*), leftmargin=*]
  \item ``Why do elements have characteristic masses?''
        was answered when the nucleus was discovered.
        What do you think determines an atom's mass?
        (Write a hypothesis now; check it in Journey~2.)

  \item ``Why do certain atoms bond together?''
        was answered by quantum mechanics.
        Based only on what you know so far, what feature
        of atoms do you think might determine their bonding?

  \item ``Why do element properties repeat periodically?''
        was Mendeleev's great mystery.
        He could describe the pattern but not explain it.
        What hidden structure might cause properties
        to cycle with increasing atomic mass?
\end{enumerate}
\end{puzzlebox}

% ─── Section: The Methods of Journey One ──────────────────────────────────────
\subsection{The Methods That Made Journey One Possible}

The content of Journey~1 matters.
But the \emph{methods} matter more,
because they are transferable to every scientific question.

\subsubsection{Method 1: Systematic Classification}

Before you can explain why things differ,
you must describe \emph{how} they differ.
Classification is not just naming.
It is the first act of pattern recognition ---
grouping things that behave similarly
and separating things that behave differently.

Lavoisier's classification of substances into elements,
compounds, and mixtures was not obvious.
For centuries, water was believed to be an element.
The classification was forced by experimental tests
(which substances resist decomposition?) not by intuition.

\subsubsection{Method 2: Quantitative Measurement}

Lavoisier's revolution was methodological:
he insisted on weighing everything,
including the gases that previous chemists had ignored.
The moment you include gases in the mass balance,
phlogiston theory collapses.

This is a general lesson:
\emph{anomalies hide in the things you have not yet measured}.
Every time chemistry advanced, it was because someone measured
something that had previously been ignored or assumed constant.

\subsubsection{Method 3: Seeking the Simplest Explanation}

Dalton did not prove atoms exist.
He showed that atoms are the simplest explanation
consistent with all the quantitative data.
This is the principle of parsimony:
among all explanations that fit the data,
prefer the simplest.

Note that ``simplest'' does not mean ``easiest to understand.''
Atoms were a stranger idea than continuous matter.
But the atomic picture required \emph{fewer arbitrary assumptions}:
once you accept discrete particles with fixed masses,
all five laws follow automatically.
Continuous matter requires separate explanations for each law.

\subsubsection{Method 4: Using Failure Productively}

Both phlogiston theory and Dalton's original atomic theory
were wrong in important ways.
But both were \emph{productively wrong}: they made precise
enough predictions that their failures could be identified
and used to build better theories.

A vague theory cannot be proven wrong.
A precise theory that makes wrong predictions
is more valuable than a vague one that makes none.

\begin{center}
\begin{tcolorbox}[
  enhanced, colback=SynthColor!5, colframe=SynthColor!60,
  width=0.85\textwidth, arc=3pt, fontupper=\sffamily
]
\centering
\textbf{The four methods of Journey One:}\\[6pt]
\small
Classify systematically $\cdot$
Measure everything $\cdot$
Seek the simplest explanation $\cdot$
Use failure productively\\[4pt]
\normalsize
These are not methods of chemistry.\\
They are methods of \emph{science}.\\
Every journey in this curriculum uses them.
\end{tcolorbox}
\end{center}

% ─── Section: What Journey One Does Not Explain ───────────────────────────────
\subsection{The Open Doors: What Journey One Cannot Tell Us}

Journey~1 ends with a picture of matter as a collection of atoms.
But the atoms themselves are treated as \emph{featureless spheres}:
they have masses and they bond together, but we know nothing
about their internal structure.

This is precisely the situation that makes Journey~2 so exciting.
Here are the open questions that atomic theory raises
but cannot answer:

\begin{openquestion}
\textbf{Question 1: What is inside an atom?}

Dalton's atoms are indivisible. But we know from electrolysis
(Davy, Faraday) that electricity can decompose compounds.
Electricity involves charged particles.
Are there charged particles inside atoms?

\vspace{0.3em}
\textit{This leads to: the discovery of the electron (Thomson, 1897).}
\end{openquestion}

\begin{openquestion}
\textbf{Question 2: How is the atom's mass distributed?}

If atoms contain electrons (which are very light),
where is the rest of the mass?
Is the atom like a uniform sphere,
or is the mass concentrated somewhere?

\vspace{0.3em}
\textit{This leads to: Rutherford's gold foil experiment (1909)
and the nuclear model.}
\end{openquestion}

\begin{openquestion}
\textbf{Question 3: Why do atoms emit light at specific colours?}

When sodium is burned, it emits a characteristic bright
yellow colour. Every element emits a unique set of colours.
Why? What is happening inside the atom?

\vspace{0.3em}
\textit{This leads to: the hydrogen spectrum, Bohr's model,
and quantum mechanics.}
\end{openquestion}

\begin{openquestion}
\textbf{Question 4: Why don't electrons spiral into the nucleus?}

Classical physics says that an electron orbiting a nucleus ---
like a planet orbiting the sun --- should continuously radiate
energy and spiral inward, destroying the atom in nanoseconds.
But atoms are stable. Why?

\vspace{0.3em}
\textit{This leads to: quantisation of energy and wave mechanics.}
\end{openquestion}

\begin{openquestion}
\textbf{Question 5: Why do element properties repeat?}

Mendeleev ordered elements by atomic mass and found that
properties recur in a pattern: every eighth element
(approximately) has similar chemistry.
Why would properties depend on mass?
What hidden structure repeats?

\vspace{0.3em}
\textit{This leads to: electron shells, orbitals,
and the full quantum model of the atom.}
\end{openquestion}

\vspace{0.5em}

Each of these questions opens an entire journey.
Journey~2 answers Questions 1, 2, 3, and 4.
Journey~3 answers Question 5.
The answers change our picture of the atom completely.

% ─── Section: Master Problem Set ─────────────────────────────────────────────
\subsection{Journey One Master Problem Set}

These problems integrate all four modules of Journey~1.
They are designed to be harder than the problems within
individual modules, and some require combining ideas
from more than one module.

\begin{problembox}[title={Integration Problems}]
\begin{enumerate}[label=\textbf{I\arabic*.}]
  \item \textbf{A mystery compound.}

        A pure white solid \textbf{X} is subjected to several tests:
        \begin{itemize}[nosep, leftmargin=2em]
          \item It dissolves completely in water to give a clear solution.
          \item Electrolysis of the solution deposits a silver metal
                at one electrode and produces a yellow-green gas at the other.
          \item The solid has a fixed melting point and fixed density,
                regardless of where it was obtained.
          \item When \SI{5.85}{\gram} of \textbf{X} is dissolved and
                the water is evaporated, exactly \SI{5.85}{\gram}
                of solid is recovered.
          \item Complete combustion analysis is not applicable
                (the compound contains no carbon or hydrogen).
        \end{itemize}

        \begin{enumerate}[label=(\alph*)]
          \item From the tests, classify \textbf{X} as element,
                compound, or mixture. Justify each step.
          \item The silver metal is identified as sodium.
                The yellow-green gas is identified as chlorine.
                What is the most likely identity of \textbf{X}?
          \item What does the electrolysis tell us about the
                bonding in \textbf{X}?
                (Hint: what must be present for electricity to
                decompose a compound?)
          \item Write a balanced equation for the electrolysis.
        \end{enumerate}

  \item \textbf{Testing the laws.}

        A student claims to have made a new compound of tin and oxygen.
        She reports the following data from three separate preparations:

        \begin{center}
        \begin{tabular}{lcc}
        \toprule
        Sample & Mass of Sn (\si{\gram}) & Mass of O (\si{\gram}) \\
        \midrule
        1 & 14.87 & 4.00 \\
        2 & 22.31 & 6.00 \\
        3 & 29.74 & 8.00 \\
        \bottomrule
        \end{tabular}
        \end{center}

        \begin{enumerate}[label=(\alph*)]
          \item Show that these data are consistent with the
                law of definite proportions.
          \item The known oxides of tin are \ce{SnO} (tin(II) oxide)
                and \ce{SnO2} (tin(IV) oxide).
                The atomic mass of Sn is \SI{118.7}{\gram\per\mol}
                and O is \SI{16.0}{\gram\per\mol}.
                Calculate the mass ratio Sn:O in each oxide.
          \item Which oxide did the student make?
          \item Another student makes \ce{SnO2}.
                If the law of multiple proportions holds,
                what mass of oxygen should combine with
                \SI{14.87}{\gram} of tin in \ce{SnO2}?
        \end{enumerate}

  \item \textbf{Gas stoichiometry.}

        Ethane (\ce{C2H6}) burns in oxygen:
        \begin{center}
        \ce{2C2H6 + 7O2 -> 4CO2 + 6H2O}
        \end{center}

        \begin{enumerate}[label=(\alph*)]
          \item What volume of \ce{O2} (at STP) is needed to
                completely burn \SI{60.0}{\gram} of ethane?
          \item What volume of \ce{CO2} (at STP) is produced?
          \item If air is 21\% \ce{O2} by volume,
                what volume of air is needed for part (a)?
          \item In practice, the reaction is carried out with
                only 85\% of the theoretical oxygen.
                What product other than \ce{CO2} and \ce{H2O}
                would you expect, and why?
                (This is why gas burners produce \ce{CO} if
                starved of air.)
        \end{enumerate}

  \item \textbf{Reasoning from first principles.}

        Before the mole was defined, chemists worked with
        ``equivalents'' --- the amount of a substance that
        reacts with a fixed reference amount of hydrogen
        or oxygen.

        Suppose you know that:
        \begin{itemize}[nosep, leftmargin=2em]
          \item \SI{8}{\gram} of oxygen reacts with \SI{1}{\gram} of hydrogen.
          \item \SI{12}{\gram} of carbon reacts with \SI{16}{\gram} of oxygen.
          \item \SI{56}{\gram} of iron reacts with \SI{24}{\gram} of oxygen.
        \end{itemize}

        And you know (from gas experiments) that:
        \begin{itemize}[nosep, leftmargin=2em]
          \item Water is \ce{H2O}.
          \item Iron oxide is \ce{Fe2O3}.
          \item Carbon dioxide is \ce{CO2}.
        \end{itemize}

        \begin{enumerate}[label=(\alph*)]
          \item Derive the atomic masses of O, C, and Fe
                relative to H = 1.
          \item Check your values against the periodic table.
          \item If you did not know the formula of iron oxide
                (i.e., you didn't know if it was \ce{FeO},
                \ce{Fe2O3}, or \ce{Fe3O4}),
                how many different atomic masses for iron
                would be consistent with the data?
          \item How does this illustrate the fundamental
                ambiguity that Avogadro's hypothesis resolved?
        \end{enumerate}
\end{enumerate}
\end{problembox}

% ─── Section: Deep problems ───────────────────────────────────────────────────
\subsection*{Deep Problems}

\begin{deepbox}
\begin{enumerate}[label=\textbf{D\arabic*.}]

  \item \textbf{The Epistemology of the Atom.}

        In 1897 the physicist Ernst Mach --- one of the most
        influential philosophers of science --- argued that
        atoms did not ``really exist'' because they could
        not be directly observed.
        He accepted that atomic theory was a useful calculating
        tool, but denied it described reality.

        Ludwig Boltzmann defended the reality of atoms his entire career.
        The debate was so bitter that it is said to have contributed
        to Boltzmann's suicide in 1906.

        \begin{enumerate}[label=(\alph*)]
          \item What criterion was Mach using for ``reality''?
                Is direct observation the only valid criterion?

          \item Dalton's atoms in 1803 were invisible and
                unmeasurable directly.
                Yet the evidence we developed in Journey~1
                was compelling.
                Make the strongest possible argument for the
                \emph{reality} of atoms using only
                pre-1900 evidence.

          \item In 1905, Einstein's analysis of Brownian motion
                gave the first way to calculate Avogadro's number
                from observable statistics of visible particles.
                Why was this considered decisive?
                What did it add that the chemical evidence lacked?

          \item Today we can image individual atoms using
                scanning tunnelling microscopes.
                Does this settle the philosophical debate?
                Or can a philosopher still argue the image
                is an inference, not a direct observation?
        \end{enumerate}

  \item \textbf{Building the First Periodic Table from Scratch.}

        Suppose you have access to all the experimental data
        available to Mendeleev in 1869:
        the atomic masses of all known elements (about 63)
        and qualitative descriptions of their chemistry.

        You are told that:
        \begin{itemize}[nosep, leftmargin=2em]
          \item Lithium, sodium, and potassium are all soft metals
                that react violently with water and form compounds
                of formula \ce{MX} (where M is the metal, X is a
                non-metal like chlorine).
          \item Fluorine, chlorine, bromine, and iodine are all
                reactive non-metals that form salts.
          \item Helium, neon, and argon form no compounds.
          \item The atomic masses increase in the order given
                within each group.
          \item Atomic masses: Li = 7, Na = 23, K = 39, F = 19,
                Cl = 35.5, Br = 80, Ar = 40.
        \end{itemize}

        \begin{enumerate}[label=(\alph*)]
          \item If you order all elements by atomic mass,
                do elements with similar properties appear
                at regular intervals?
                Calculate the mass gaps within each group above.

          \item What pattern do you see in the gaps?

          \item If the pattern continues, predict the atomic mass
                of an undiscovered element that should appear
                between calcium (40) and titanium (48).
                (This was Mendeleev's prediction of ``eka-calcium,''
                later discovered as scandium with mass 45.)

          \item Argon (mass 40) comes \emph{before} potassium (mass 39)
                by chemical behaviour, but \emph{after} potassium
                by atomic mass.
                This was a genuine anomaly for Mendeleev.
                Propose a hypothesis for why atomic mass order
                might sometimes differ from chemical order.
                (Hint: what do we know about isotopes?)
        \end{enumerate}

  \item \textbf{Counting the Invisible: A Thought Experiment.}

        You are a chemist in 1860.
        You have the following data:
        \begin{itemize}[nosep, leftmargin=2em]
          \item The density of hydrogen gas at STP:
                \SI{0.0899}{\gram\per\litre}.
          \item The fact that hydrogen is diatomic (\ce{H2}).
          \item The atomic mass of hydrogen: 1.008 u.
          \item The measured size of a hydrogen atom from
                kinetic theory: approximately
                $1.2 \times 10^{-10}\ \si{\metre}$ diameter.
        \end{itemize}

        \begin{enumerate}[label=(\alph*)]
          \item Using only the density and the fact that
                1 mole of gas occupies \SI{22.4}{\litre} at STP,
                calculate the mass of one mole of \ce{H2}.
                From this, deduce the mass of one hydrogen atom in grams.

          \item If a hydrogen atom has diameter
                $d = 1.2\times10^{-10}$ m,
                estimate its volume (treat it as a sphere).

          \item From the mass and volume of one atom,
                estimate the density of ``solid hydrogen''
                if atoms were packed with no gaps.
                How does this compare with the actual density
                of liquid hydrogen ($\approx \SI{71}{\gram\per\litre}$)?
                What does the discrepancy tell you?

          \item From parts (a) and the molar mass,
                calculate Avogadro's number.
                How close is your answer to the modern value?
        \end{enumerate}

\end{enumerate}
\end{deepbox}

% ─── Section: Journey One Synthesis ──────────────────────────────────────────
\subsection*{Journey One: The Complete Synthesis}

\begin{synthbox}
\textbf{The single sentence of Journey One:}
\begin{center}
\large\sffamily\bfseries
Matter is made of discrete, indivisible atoms;
each element has one type;\\
compounds are fixed combinations;
and all of this was forced upon us\\
by measurement, not imagination.
\end{center}

\vspace{0.6em}
\textbf{The five laws, one by one:}

\vspace{0.3em}
\begin{tabular}{lp{9cm}}
\toprule
\textbf{Law} & \textbf{Why it demands atoms} \\
\midrule
Conservation of mass &
  If atoms rearrange but are never destroyed, mass is always conserved.
  No other simple model predicts this for all reactions. \\[5pt]
Definite proportions &
  If a compound is always the same atom ratio, its mass ratio is fixed.
  Continuous matter has no reason to produce fixed ratios. \\[5pt]
Multiple proportions &
  If you must use whole numbers of atoms, ratios between
  related compounds must be simple integers.
  Irrational ratios are impossible. \\[5pt]
Combining volumes &
  If equal volumes contain equal numbers of molecules,
  volume ratios directly reflect molecule ratios.
  This only works if molecules are discrete. \\[5pt]
Avogadro's hypothesis &
  Resolves the diatomic-molecule problem and fixes all
  molecular formulas. Makes the mole scale self-consistent. \\
\bottomrule
\end{tabular}

\vspace{0.6em}
\textbf{The open questions that lead us forward:}
\begin{itemize}[nosep, leftmargin=2em]
  \item What is inside an atom? $\longrightarrow$ Journey~2
  \item Why do properties repeat? $\longrightarrow$ Journey~3
  \item Why do atoms bond? $\longrightarrow$ Journey~4
  \item What drives reactions? $\longrightarrow$ Journey~5
  \item What can carbon do? $\longrightarrow$ Journeys~6 and~7
  \item What can we build? $\longrightarrow$ Journey~8
\end{itemize}
\end{synthbox}

% ─── Final Reflection ─────────────────────────────────────────────────────────
\subsection*{Final Reflection for Journey One}

\begin{reflectionbox}
These questions are meant to be hard.
Take time. Write at length.
There are no short correct answers.

\begin{enumerate}[label=\textbf{R\arabic*.}]
  \item \textbf{The logic of the atomic hypothesis.}

        Dalton did not discover atoms.
        He showed that atoms were the \emph{simplest explanation}
        of the data.
        Write a careful argument for why the five laws together
        constitute strong evidence for atoms,
        even though no individual law proves atoms exist.
        Use the word ``convergence.''

  \item \textbf{What makes a theory scientific?}

        Phlogiston theory explained a great deal of chemistry
        and was believed by excellent scientists.
        Dalton's atomic theory also explained a great deal
        and was partially wrong (water as \ce{HO}).

        What distinguishes a \emph{good scientific theory}
        from a \emph{bad one}?
        Is it just whether the theory is correct?
        Or is there something else --- a property that makes
        a theory worth holding and worth testing?

  \item \textbf{The role of instrumentation.}

        Many of the advances in Journey~1 required new instruments:
        precision balances (Lavoisier), gas collection apparatus
        (Priestley, Cavendish), electrolysis equipment (Davy).

        Write a general principle about the relationship between
        instruments and scientific discovery.
        What determines when a new instrument is \emph{needed}?
        What kinds of discoveries are impossible without better tools?

  \item \textbf{Looking forward.}

        Based on what you know from Journey~1,
        write your best current hypothesis for why
        elements have the chemical properties they do.
        Be specific: what feature of an atom do you think
        determines whether it is reactive or inert,
        whether it bonds to many other atoms or few?

        You do not need to be correct.
        You need to be \emph{precise enough that your hypothesis
        makes a prediction that could be tested}.
        Write the test.
\end{enumerate}
\end{reflectionbox}

% ─── Closing page ─────────────────────────────────────────────────────────────
\vspace{2em}
\begin{center}
\begin{tikzpicture}
  \node[
    draw=JourneyColor, fill=JourneyColor, rounded corners=8pt,
    minimum width=13cm, minimum height=5.5cm
  ] (bg) at (0,0) {};

  \node[
    font=\sffamily\large\bfseries, text=white,
    text width=12cm, align=center
  ] at (0, 1.4) {Journey One is complete.};

  \node[
    font=\sffamily\small, text=white!80,
    text width=11.5cm, align=center
  ] at (0, 0.2) {
    You began with the observation that substances differ.\\
    You end knowing that matter is discrete, that atoms are real,\\
    and that the evidence for this --- accumulated over two centuries\\
    by dozens of scientists in different countries ---\\
    converges on a single picture from five independent directions.
  };

  \node[
    font=\sffamily\normalsize\bfseries\itshape, text=DiscoveryColor!80,
    text width=11.5cm, align=center
  ] at (0, -1.5) {
    The atoms are real.\\
    But we do not yet know what they look like inside.\\[5pt]
    That is the question of Journey Two.
  };
\end{tikzpicture}
\end{center}

\end{document}